% P05.1 strict-vertex integral-slope nonregularity proof, 26 September 2026.

\begin{proof}[\Cref{thm:frontier-polyhedral-strict-vertex-integral-nonregularity}]
The balanced weight vector belongs to $K_v\setminus\{0\}$, so the slice
\[
              B:=\left\{\lambda\in K_v:\sum_e\lambda_e=1\right\}
\]
is nonempty and compact.  By
\eqref{eq:frontier-polyhedral-strict-vertex}, the continuous function
$\lambda\mapsto\sum_e\lambda_es_e$ has a positive minimum $\eta$ on
$B$.  Every integer $N\geq\max\{2,\eta^{-1}\}$ therefore satisfies
\eqref{eq:frontier-polyhedral-strict-vertex-choice}.

The function $u_{v,R,N}$ is continuous: all incident germs have the common
value $c$ at $v$, and its value is zero at every $x_e(R)$.  Its only slopes
are $0$, $-1$, and $0$, and all its breakpoints are rational.  Hence it is a
bounded literal piecewise-affine function.

Denote the right-hand side of
\eqref{eq:frontier-polyhedral-strict-vertex-envelope} by $P$.  Its outgoing
slope at $v$ along $e$ is $s_e-1/N$, so
\eqref{eq:frontier-polyhedral-strict-vertex-choice} and the local slope
criterion of
\Cref{lem:frontier-polyhedral-one-vertex-convexity} give polyhedral
convexity at $v$.  At every $x_e(R)$ the slope increases from
$s_e-1/N$ to $s_e$; all remaining local inequalities are those of
$\gamma$.  Thus $P$ is polyhedrally convex and differs from $\gamma$ by a
bounded function.

We next check the obstacle inequality.  On the constant part of each star
segment,
\[
                         \frac{R-t}{N}\leq\frac RN=c.
\]
On its final part, $(R-t)/N\leq R-t$, and off the star equality holds.
Consequently $P\leq\gamma+u_{v,R,N}$, so $P$ is an admissible envelope
competitor.

Conversely, let $h\leq\gamma+u_{v,R,N}$ be any admissible polyhedrally
convex competitor.  At the two endpoints of every star segment,
\[
              h(v)\leq P(v),\qquad h(x_e(R))\leq P(x_e(R)).
\]
Ordinary convexity on $x_e([0,R])$ bounds $h$ by the affine chord joining
these endpoint bounds.  Since $\gamma$ is affine there, this chord is
exactly $P$.  Off the closed star the obstacle itself equals $P$.  Hence
$h\leq P$ everywhere, proving
\eqref{eq:frontier-polyhedral-strict-vertex-envelope}.

It remains to prove that this envelope is not literally regularizable.  Fix
one incident germ $e$, put $q:=s_e-1/N$, and abbreviate
\[
 \begin{aligned}
 e_0&:=\psi(v)-P(v),\\
 e_c&:=\psi(x_e(c))-P(x_e(c)),\\
 e_R&:=\psi(x_e(R))-P(x_e(R)).
 \end{aligned}
\]
The left derivative
$m:=(\psi\circ x_e)'_-(c)$ is an integer.  If $m\leq s_e-1$, convexity on
$[0,c]$ gives
\[
 e_c-e_0
 \leq (m-q)c
 \leq-\frac{N-1}{N}\frac RN
 =-\frac{(N-1)R}{N^2}.
\]
If $m\geq s_e$, monotonicity of the one-sided derivatives gives
\[
 e_R-e_c
 \geq(m-q)(R-c)
 \geq\frac1N\frac{(N-1)R}{N}
 =\frac{(N-1)R}{N^2}.
\]
The two cases exhaust the integral values of $m$.  In either case, at least
one of the errors occurring in the relevant difference has absolute value
at least $(N-1)R/(2N^2)$.  This proves
\eqref{eq:frontier-polyhedral-strict-vertex-separation}.  No sequence of
literal integral-slope convex piecewise-affine functions can converge
uniformly to the envelope, so literal regularity fails.
\end{proof}

\begin{proof}[\Cref{cor:frontier-polyhedral-rank-one-graph-integral-nonregularity}]
Refine the finite structure so that $\gamma$ is affine on every edge.  A
new bivalent vertex has two opposite primitive directions and equal
inherited balancing weights, so its local relation cone remains
one-dimensional.  The positive measure
$\MA_{\mathrm{poly}}(\gamma)$ is atomic at the vertices and has total mass
$\deg(\gamma)>0$.  Choose a vertex $v$ of positive mass, and write $s_e$
for the outgoing slopes of $\gamma$ there.  Since
$K_v=\mathbf R_{\geq0}w_v$, every nonzero $\lambda\in K_v$ has the form
$\lambda=aw_v$ with $a>0$, and
\[
                  \sum_e\lambda_es_e
                  =a\,\MA_{\mathrm{poly}}(\gamma)\{v\}>0.
\]
The hypothesis of
\Cref{thm:frontier-polyhedral-strict-vertex-integral-nonregularity} holds
at $v$, so literal regularity fails.
\end{proof}
